\section{Data Structure}
\subsection{Sparse Table}
Max/Min/GCD Query
\begin{minted}{c++}
struct SparseTable {
  SparseTable(vector<int>& data)
  {
    int n = data.size();
    int maxLog = log2(n) + 1;
    table.assign(n, vector<int>(maxLog));
    for (int i = 0; i < n; ++i) {
      table[i][0] = data[i];
    }
    for (int j = 1; (1 << j) <= n; ++j) {
      for (int i = 0; (i + (1 << j) - 1) < n; ++i) {
        table[i][j] = max(table[i][j - 1], table[i +
        (1 << (j - 1))][j - 1]);
      }
    }
  }
  int query(int L, int R) {
    int j = log2(R - L + 1);
    return max(table[L][j], table[R - (1 << j) +
    1][j]);
  }
  vector<vector<int>> table;
};
\end{minted}
\subsection{Sparse Table 2D}
\begin{minted}{c++}
struct SparseTable2D {
  static const int LG = 10;
  int n, m;
  vector<vector<int>> a;
  vector<int> lg2;
  vector<vector<vector<vector<int>>>> st;
  SparseTable2D(int n, int m) : n(n), m(m) {
    a.assign(n, vector<int>(m, 0));
    lg2.resize(max(n, m) + 1, 0);
    for (int i = 2; i <= max(n, m); i++) lg2[i] =
    lg2[i >> 1] + 1;
    st.assign(n, vector<vector<vector<int>>>(m,
    vector<vector<int>>(LG, vector<int>(LG,
    0))));
  }
  void set_val(int i, int j, int val) { a[i][j] = val; }
  void build() {
    for (int i = 0; i < n; i++)
      for (int j = 0; j < m; j++)
        st[i][j][0][0] = a[i][j];
    for (int a = 0; a < LG; a++) {
      for (int b = 0; b < LG; b++) {
        if (a + b == 0) continue;
        for (int i = 0; i + (1 << a) <= n; i++) {
          for (int j = 0; j + (1 << b) <= m; j++) {
            if (!a)
              st[i][j][a][b] = max(st[i][j][a][b-1],
              st[i][j + (1 << (b - 1))][a][b - 1]);
            else
              st[i][j][a][b] = max(st[i][j][a-1][b],
              st[i + (1 << (a - 1))][j][a - 1][b]);
          }
        }
      }
    }
  }
  int query(int x1, int y1, int x2, int y2) {
    int a = lg2[x2 - x1 + 1], b = lg2[y2 - y1 + 1];
    int m1 = max(st[x1][y1][a][b], st[x2 - (1 << a)
    + 1][y1][a][b]);
    int m2 = max(st[x1][y2 - (1 << b) + 1][a][b],
    st[x2 - (1 << a) + 1][y2 - (1 << b) + 1][a][b]);
    return max(m1, m2);
  }
};
// st.set_val(i, j, grid[i][j]);
// st.build(); // st.query(0, 0, 1, 2); // max in
// submatrix [(0,0) to (1,2)]
\end{minted}
\subsection{Merge Sort Tree (Values $\le$ K in Range)}
\begin{minted}{c++}
struct MergeSortTree {
  int n;
  vector<vector<int>> tree;
  MergeSortTree(const vector<int>& arr) {
    n = arr.size();
    tree.resize(4 * n);
    build(1, 0, n - 1, arr);
  }
  void build(int node, int l, int r, const vector<int>& arr) {
    if (l == r) {
      tree[node].push_back(arr[l]);
      return;
    }
    int mid = (l + r) / 2;
    build(2 * node, l, mid, arr);
    build(2 * node + 1, mid + 1, r, arr);
    merge(tree[2 * node].begin(), tree[2 * node].end(),
          tree[2 * node + 1].begin(), tree[2 * node + 1].end(),
          back_inserter(tree[node]));
  }
  int query(int node, int l, int r, int ql, int qr, int k) {
    if (l > qr || r < ql) return 0;
    if (l >= ql && r <= qr) {
      return upper_bound(tree[node].begin(),
      tree[node].end(), k) - tree[node].begin();
    }
    int mid = (l + r) / 2;
    return query(2 * node, l, mid, ql, qr, k) +
           query(2 * node + 1, mid + 1, r, ql, qr, k);
  }
};
\end{minted}
\subsection{K-th Smallest in Range (Index SegTree)}
\begin{minted}{c++}
struct KthSmallest {
  int n;
  vector<vector<int>> seg;
  KthSmallest(const vector<int>& arr) {
    n = arr.size();
    seg.resize(4 * n);
    vector<pair<int, int>> B(n);
    for (int i = 0; i < n; i++) B[i] = {arr[i], i};
    sort(B.begin(), B.end());
    build(0, 0, n - 1, B);
  }
  void build(int node, int l, int r,
  const vector<pair<int, int>>& B) {
    if (l == r) {
      seg[node].push_back(B[l].second);
      return;
    }
    int mid = (l + r) / 2;
    build(2 * node + 1, l, mid, B);
    build(2 * node + 2, mid + 1, r, B);
    merge(seg[2 * node + 1].begin(),
    seg[2 * node + 1].end(),
          seg[2 * node + 2].begin(),
          seg[2 * node + 2].end(),
          back_inserter(seg[node]));
  }
  int query(int node, int st, int end, int l, int r,
  int k) {
    if (st == end) return seg[node][0];
    int p = upper_bound(seg[2 * node + 1].begin(),
    seg[2 * node + 1].end(), r)
          - lower_bound(seg[2 * node + 1].begin(),
          seg[2 * node + 1].end(), l);
    int mid = (st + end) / 2;
    if (p >= k)
      return query(2 * node + 1, st, mid, l, r, k);
    else return query(2 * node + 2, mid + 1, end, l,
    r, k - p);
  }
};
\end{minted}
\subsection{Persistent Segment Tree}
\begin{minted}{c++}
struct PST {
  struct Node { int l, r, val; };
  vector<Node> t;
  int T = 0;
  PST(int max_nodes) { t.resize(max_nodes); }
  int build(int b, int e) {
    int cur = T++;
    if (b == e) return cur;
    int m = (b + e) / 2;
    t[cur].l = build(b, m);
    t[cur].r = build(m + 1, e);
    return cur;
  }
  int update(int pre, int b, int e, int i, int v) {
    int cur = T++;
    t[cur] = t[pre];
    if (b == e) { t[cur].val += v; return cur; }
    int m = (b + e) / 2;
    if (i <= m) t[cur].l = update(t[pre].l, b, m, i, v);
    else t[cur].r = update(t[pre].r, m + 1, e, i, v);
    t[cur].val = t[t[cur].l].val + t[t[cur].r].val;
    return cur;
  }
  int query(int pre, int cur, int b, int e, int k) {
    if (b == e) return b;
    int cnt = t[t[cur].l].val - t[t[pre].l].val;
    int m = (b + e) / 2;
    if (cnt >= k)
      return query(t[pre].l, t[cur].l, b, m, k);
    return query(t[pre].r, t[cur].r, m + 1, e,
    k - cnt);
  }
};
\end{minted}
\subsection{Fenwick Tree}
Point update and range query.
\begin{minted}{c++}
template<typename T>
struct BIT {
  int n;
  vector<T> bit;
  BIT(int n = 0) { init(n); }
  void init(int n_) { n = n_; bit.assign(n + 1, 0); }
  void add(int i, T v) {
    for (; i <= n; i += i & -i) bit[i] += v; }
  T sum(int i) {
    T r = 0;
    for (; i > 0; i -= i & -i) r += bit[i];
    return r;
  }
  T sum(int l, int r) { return sum(r) - sum(l - 1); }
};
\end{minted}
Usage: \texttt{BIT<ll> fenw(5);}
\subsection{Fenwick 2D}
\begin{minted}{c++}
struct BIT2D {
  int n, m;
  vector<vector<ll>> bit;
  BIT2D(int n, int m): n(n), m(m) {
    bit.assign(n + 1, vector<ll>(m + 1, 0));
  }
  void add(int x, int y, ll val) {
    for (int i = x; i <= n; i += i & -i)
      for (int j = y; j <= m; j += j & -j)
        bit[i][j] += val;
  }
  ll sum(int x, int y) {
    ll res = 0;
    for (int i = x; i > 0; i -= i & -i)
      for (int j = y; j > 0; j -= j & -j)
        res += bit[i][j];
    return res;
  }
  ll query(int x1, int y1, int x2, int y2) {
    return sum(x2, y2) - sum(x1 - 1, y2) -
    sum(x2, y1 - 1) + sum(x1 - 1, y1 - 1);
  }
};
BIT2D fenw2(4, 4);
fenw2.add(1, 1, 5);
fenw2.add(3, 2, 7);
cout << fenw2.query(1, 1, 3, 3) << "\n";
\end{minted}
\subsection{BIT Range Update \texorpdfstring{\&}{&} Range Query}
\begin{minted}{c++}
using ll = long long;
ll bit1[200005], bit2[200005], n, m;
void update(ll bit[], ll i, ll val) {
  for (; i <= n; i += (i & (-i))) bit[i] += val;
}
void range_update(ll l, ll r, ll val) {
  update(bit1, l, val);
  update(bit1, r + 1, -val);
  update(bit2, l, val * (l - 1));
  update(bit2, r + 1, -(r * val));
}
ll sum(ll bit[], ll i) {
  ll re = 0;
  for (; i > 0; i -= (i & (-i))) re += bit[i];
  return re;
}
ll prefix_sum(ll i) {
  return sum(bit1, i) * i - sum(bit2, i);
}
ll range_sum(ll l, ll r) {
  return prefix_sum(r) - prefix_sum(l - 1);
}
\end{minted}
\texttt{range\_update(l,r,val);}\\
\texttt{range\_sum(l,r);}
\subsection{Implicit Treap}
\begin{minted}{c++}
mt19937
rnd(chrono::steady_clock::now().time_since_epoch
().count());
struct Node {
  int64_t prior, data, sum, toProp;
  int size; bool rev; Node *l, *r;
  Node(int64_t val) {
    prior = rnd(), data = sum = val;
    size = 1; rev = toProp = 0;
    l = r = nullptr;
  }
};
struct Treap {
  Node *root;
  Treap() : root(nullptr) {}
  int size(Node *&t) { return t ? t->size : 0; }
  int64_t sum(Node *&t) { return t ? t->sum : 0; }
  void update(Node *&t) {
    if (!t) return;
    t->size = size(t->l) + size(t->r) + 1;
    t->sum = t->data + t->size * t->toProp;
    if (t->l) t->sum += t->l->sum +
    t->l->size * t->l->toProp;
    if (t->r) t->sum += t->r->sum +
    t->r->size * t->r->toProp;
  }
  void push(Node *&t) {
    if (!t) return;
    if (t->rev) {
      t->rev = false; swap(t->l, t->r);
      if (t->l) t->l->rev ^= true;
      if (t->r) t->r->rev ^= true;
    }
    if (t->toProp) {
      t->data += t->toProp;
      if (t->l) t->l->toProp += t->toProp;
      if (t->r) t->r->toProp += t->toProp;
      t->toProp = 0;
    }
  }
  void split(Node *t, Node *&l, Node *&r, int key,
  int add = 0) {
    if (!t) l = r = nullptr;
    else {
      push(t);
      if (key <= add + size(t->l))
        split(t->l, l, t->l, key, add), r = t;
      else split(t->r, t->r, r, key,
      add + 1 + size(t->l)), l = t;
      update(t);
    }
  }
  void merge(Node *&t, Node *l, Node *r) {
    push(l), push(r);
    if (!l || !r) t = l ? l : r;
    else if (l->prior > r->prior)
      merge(l->r, l->r, r), t = l;
    else merge(r->l, l, r->l), t = r;
    update(t);
  }
  void insert(Node *&t, int key, int64_t val) {
    Node *left, *right;
    split(t, left, right, key);
    merge(left, left, new Node(val));
    merge(t, left, right);
  }
  void reverse(Node *&t, int l, int r) {
    Node *left, *mid, *right;
    split(t, left, right, l);
    split(right, mid, right, r - l + 1);
    assert(mid); mid->rev ^= true;
    merge(t, left, mid);
    merge(t, t, right);
  }
  int erase(Node *&t, int key) {
    assert(size(t) > key);
    Node *left, *mid, *right;
    split(t, left, right, key);
    split(right, mid, right, 1);
    merge(t, left, right);
    int ret = mid->data;
    delete mid; return ret;
  }
  int64_t query_sum(Node *&t, int l, int r) {
    Node *left, *mid, *right;
    split(t, left, right, l);
    split(right, mid, right, r - l + 1);
    int64_t ret = mid->sum;
    merge(t, left, mid);
    merge(t, t, right);
    return ret;
  }
  void rangeAdd(Node *&t, int l, int r, int64_t val) {
    Node *left, *mid, *right;
    split(t, left, right, l);
    split(right, mid, right, r - l + 1);
    mid->toProp += val;
    merge(t, left, mid);
    merge(t, t, right);
  }
  void set(Node *&t, int p, int64_t val) {
    Node *left, *mid, *right;
    split(t, left, right, p);
    split(right, mid, right, 1);
    merge(t, left, new Node(val));
    merge(t, t, right);
  }
};
\end{minted}
\subsection{Mo's Algorithm}
\begin{minted}{c++}
sort(cu.begin(), cu.end(), [&](auto &x, auto &y){
  if (x.l / B != y.l / B)
    return x.l / B < y.l / B;
  return (x.l / B) % 2 ? x.r > y.r : x.r < y.r;
});
int R = -1, L = 0;
while (R < r) add(++R);
while (L > l) add(--L);
while (R > r) rem(R--);
while (L < l) rem(L++);
\end{minted}
